\documentclass[10pt,twocolumn]{article}
\usepackage[T1]{fontenc}
\usepackage[utf8]{inputenc}
\usepackage{lmodern}
\usepackage[margin=1.8cm,top=2.4cm,bottom=2cm,columnsep=0.6cm]{geometry}
\usepackage{fancyhdr}
\usepackage{titlesec}
\usepackage{booktabs}
\usepackage{amsmath,amssymb,amsfonts}
\usepackage{microtype}
\usepackage{xcolor}
\usepackage{mdframed}
\usepackage{parskip}
\usepackage{enumitem}
\usepackage{hyperref}

\definecolor{rulered}{RGB}{180,30,30}
\definecolor{boxgray}{RGB}{245,245,245}
\definecolor{keywordgray}{RGB}{100,100,100}

\pagestyle{fancy}
\fancyhf{}
\fancyhead[L]{\textsc{\small Claude Mag}}
\fancyhead[C]{\small\textit{Machine Learning Research Digest}}
\fancyhead[R]{\small 2026-09-07}
\fancyfoot[C]{\small\thepage}
\renewcommand{\headrulewidth}{0.8pt}

\titleformat{\section}{\normalsize\bfseries\color{rulered}}{}{0pt}{}%
  [\vspace{-4pt}\color{rulered}\rule{\columnwidth}{0.4pt}]
\titlespacing{\section}{0pt}{8pt}{4pt}

\hypersetup{colorlinks=true,urlcolor=rulered,linkcolor=black}

\begin{document}
\setlength{\parindent}{0pt}
\setlength{\parskip}{4pt}

{\small\color{keywordgray}\textbf{Keywords:}
  on-policy distillation \textbullet{}
  knowledge distillation \textbullet{}
  LLM training efficiency \textbullet{}
  state coverage}\par\vspace{4pt}

{\LARGE\bfseries\raggedright
  Rethinking On-Policy Distillation:\\
  One Training Example}\par\vspace{4pt}

{\large\itshape\raggedright
  A Single Query Recovers Most of Full-Dataset
  OPD Gain --- Revealing That LLM Distillation
  Is Algorithm-Starved, Not Data-Starved}\par\vspace{6pt}

{\small\textbf{By} Zixuan Fu, Bingxiang He, Yuxin Zuo, Haohuan Huang,
  Jinqian Zhang, Ruhang Xiao, Cheng Qian, Qinyu Luo, Huan-ang Gao,
  Yudong Wang, Zhiyuan Liu, Ning Ding, Chaojun Xiao
  (Tsinghua University et al.)}\par\vspace{2pt}

{\small\color{keywordgray}arXiv:2609.04172 \textbullet{} Preprint, September 2026}
\par\vspace{2pt}\noindent{\color{rulered}\rule{\columnwidth}{0.6pt}}\vspace{6pt}

On-policy distillation (OPD) of large language models is widely believed
to require large, diverse training datasets. This paper refutes that
assumption: training OPD on a single query recovers most of the accuracy
gain achieved by full-dataset training on mathematical reasoning, code,
and general benchmarks across multiple model families. The finding reveals
that OPD is \emph{data-overfed but algorithm-starved} --- the bottleneck
lies in the optimization procedure, not the volume of training queries.

\begin{mdframed}[backgroundcolor=boxgray,linecolor=rulered,linewidth=1pt,
    innerleftmargin=6pt,innerrightmargin=6pt,
    innertopmargin=6pt,innerbottommargin=6pt]
\textbf{\small AT A GLANCE}\par\vspace{4pt}
\begin{description}[leftmargin=0pt,labelsep=4pt,itemsep=2pt,parsep=0pt,
    font=\small\bfseries]
  \item[Problem:] OPD is assumed to require large diverse datasets;
    no theory or experiment had tested the data-minimal limit.
  \item[Why it matters:] If one query suffices, OPD can be applied
    at very low data cost, democratizing teacher-guided alignment.
  \item[Method:] Fix training steps; vary query count from 1 to
    the full dataset; measure accuracy and state coverage.
  \item[Result:] One query achieves 71.5\% state coverage and most of
    full-data OPD's gain; 16 queries reach 98.9\% and match full data.
  \item[Evidence:] Experiments across MATH-500, GSM8K, LiveCodeBench,
    MMLU-Pro with Qwen2.5-7B, Llama-3.1-8B, DeepSeek-R1-Distill-7B.
\end{description}
\end{mdframed}

\section{The Big Picture}

On-policy distillation combines student-generated rollouts with dense
token-level supervision from a teacher LLM. The student generates
responses, the teacher evaluates them, and the student updates its
parameters using those signals --- with rollouts regenerated at each
iteration from the updated student.

The conventional wisdom is that OPD needs a large and diverse set of
training queries to cover the space of possible inputs and to produce
varied rollouts for learning from. This view motivates the standard
practice of assembling large question banks for distillation.

This paper asks a simple question: what happens at the data-minimal
limit, i.e.\ when OPD trains on exactly one query? The answer, that
one-shot OPD keeps improving for hundreds of steps and recovers most
of full-data OPD's gain, reframes the problem of LLM alignment
distillation entirely.

\section{Why Existing Approaches Fall Short}

Current OPD practice assembles datasets of thousands of training queries
(e.g.\ 512--4096 mathematical problems), under the assumption that
broad coverage is necessary for generalization. This creates two problems:

First, large datasets impose \emph{data collection costs}: curating
high-quality, diverse questions for every new target domain requires
significant effort.

Second, large datasets may \emph{mask algorithmic limitations}: if the
training data is the bottleneck, then improving the algorithm yields
little gain unless data is first increased. This paper reveals that
the data is \emph{not} the bottleneck, so algorithmic improvements
(better teachers, better update rules, better rollout strategies) are
the correct direction for future work.

\section{Core Mechanism}

The key insight is that even with a fixed single query $q^*$, the
rollout distribution $\pi_\theta(\cdot\mid q^*)$ changes at every
gradient step as $\theta$ is updated. This \emph{dynamic rollout
diversity} provides gradient signal across different student behaviours
even when the input prompt is fixed.

The authors measure this via \textbf{state coverage}: the fraction of
intermediate rollout states visited by a given query set that are also
visited by the full dataset:
\[
  \mathrm{Coverage}(Q) = \frac{|S(Q) \cap S(Q_\mathrm{full})|}{|S(Q_\mathrm{full})|}
\]
One query achieves $\mathrm{Coverage} = 71.5\%$ --- most of the state
space is reachable from a single starting point as the model evolves.

\section{Technical Formulation}

The OPD training objective is a token-level imitation loss:
\[
  \mathcal{L}_\mathrm{OPD}(\theta) = -\mathbb{E}_{q \sim \mathcal{D}}\,
  \mathbb{E}_{y \sim \pi_\theta(\cdot|q)}
  \sum_t \log \pi_\mathrm{teacher}(y_t \mid q, y_{<t})
\]
where $\pi_\theta$ is the student, $\pi_\mathrm{teacher}$ is the frozen
teacher, and $y$ is a student-generated rollout.

At the single-query limit with query $q^*$, the gradient at step $k$ is:
\[
  \nabla_\theta \mathcal{L} =
  -\mathbb{E}_{y \sim \pi_\theta(\cdot|q^*)}
  \sum_t \nabla_\theta \log \pi_\theta(y_t \mid q^*, y_{<t})
  \cdot \log \pi_\mathrm{teacher}(y_t \mid q^*, y_{<t})
\]
Despite the fixed $q^*$, the expectation is over the current student
rollout distribution, which shifts every step --- providing effective
gradient diversity.

\subsection{Coverage Saturation}

Empirically, coverage $C(N)$ as a function of query count $N$ follows
an approximate power law:
\[
  C(N) \approx 1 - \alpha N^{-\beta}
\]
with fitted constants $\alpha, \beta$. Empirically, $C(1) \approx 0.715$
and $C(16) \approx 0.989$.

\section{Learning or Inference Procedure}

The authors fix the total training budget (number of gradient steps)
and vary only the number of unique training queries. For each query count
$N \in \{1, 4, 16, 64, 512\}$:
\begin{enumerate}[itemsep=2pt,parsep=0pt]
  \item Uniformly sample (or select by diversity) $N$ queries from the
    training pool.
  \item Run OPD for the fixed number of steps on those $N$ queries.
  \item Evaluate on held-out benchmarks (MATH-500, GSM8K, LiveCodeBench,
    MMLU-Pro).
  \item Measure state coverage by intersecting rollout state sets.
\end{enumerate}

Query selection at $N > 1$ uses a semantic diversity criterion (maximum
marginal relevance on embedding space) to maximise coverage gain.

\section{What the Guarantee Says}

The paper does not prove formal guarantees but establishes an empirical
regularity: OPD accuracy as a function of query count follows the same
coverage power law. Because accuracy and coverage co-move, maximising
\emph{semantic diversity} of a small query set --- rather than maximising
dataset size --- is the correct strategy for efficient OPD.

The corollary is that OPD improvements will be unlocked by better
algorithms (improved rollout strategies, better acceptance rules,
curriculum over query difficulty) rather than by larger datasets.

\section{Experimental Findings}

\textbf{Accuracy vs. query count (Qwen2.5-7B, MATH-500):}

\begin{center}
\begin{tabular}{rcc}
\toprule
Queries & MATH-500 & Coverage \\
\midrule
1       & 68.4\%   & 71.5\% \\
4       & 70.1\%   & 89.2\% \\
16      & 71.8\%   & 98.9\% \\
512     & 72.3\%   & 100\%  \\
\bottomrule
\end{tabular}
\end{center}

The accuracy gain from 1 query to 512 queries is only 3.9 percentage
points, while the gain from the pre-OPD baseline to 1-query OPD is
substantially larger --- confirming that most of the alignment benefit
comes from the OPD mechanism itself, not the data diversity.

Results are consistent across Llama-3.1-8B and DeepSeek-R1-Distill-7B
student models, and across GSM8K, LiveCodeBench, and MMLU-Pro benchmarks.

\section{Ablations and Interpretation}

\textbf{Random vs.\ semantic query selection:} At $N=16$, randomly
sampled queries underperform semantically diverse queries by 1.7
percentage points (MATH-500). Coverage is the active ingredient.

\textbf{Step budget vs.\ data size:} One-shot OPD requires $\sim$2$\times$
more gradient steps to match full-data OPD accuracy at $N=16$, but
converges to the same value --- showing the algorithm can compensate for
data sparsity given sufficient compute.

\textbf{Teacher quality:} Weaker teachers (same-size models) yield
diminished but still positive one-shot gains. The mechanism does not
require a very large teacher.

\textbf{Alignment pace:} Plotting validation accuracy vs.\ steps for
$N=1$ and $N=512$, both curves plateau at the same rate --- confirming
the algorithm bottleneck hypothesis.

\section*{Reference}

Zixuan Fu, Bingxiang He, Yuxin Zuo et al.
\textbf{Rethinking On-Policy Distillation of Large Language Models II:
One Training Example.}
arXiv:2609.04172, September 2026.
\url{https://arxiv.org/abs/2609.04172}

\end{document}
